\chapter[Introduction]{Introduction}\label{ch:sec:introduction}

A side-channel analyst measures the power consumption of a chip and
obtains a distribution over possible secret keys.  A location-privacy
researcher reports a noisy GPS coordinate and wants to guarantee that
an adversary cannot distinguish nearby locations.  A cryptographer
designs an encryption scheme and needs to know whether any ciphertext
leaks information about the plaintext.

All three problems share the same mathematical core: given an encoder
that produces outputs from secret inputs, how far is the resulting
distribution from being independent of the input?  This book develops
a unified geometric framework that answers this question for any
deterministic encoder.

The answer is the \emph{Universal Optimality Defect} (UOD), a single
number that measures the distance from a distribution to the ideal of
perfect secrecy.  The framework that produces this number---weighted
incidence structures, their forced Hilbert-space geometry, and the
canonical operator calculus---is the subject of this book.

\section{Main Contributions}\label{ch:sec:main-contributions}

The book makes twelve principal contributions.

\paragraph{1. Complete classification of universally optimal encoders.}
Before this work, only one family of universally optimal deterministic
encoders was known---the Apollonian Cell Encoders (ACE)
\cite{biondi2018apollonian}.  The present book closes the problem.  A
weighted incidence structure is realised by a deterministic encoder if
and only if its induced multiplicities are integers and its weights
satisfy the linear balance equation $B\beta = Kw$ on the incidence
graph (Theorem~\ref{ch:thm:realizability}).  This is a complete,
checkable characterisation.

\paragraph{2. The WIS as an intrinsic invariant.}
The WIS is not a modelling convenience.  It emerges inevitably from the
probabilistic data of any universally optimal encoder.  The
representation theorem (Theorem~\ref{ch:thm:representation}) shows that
every such encoder induces a unique WIS, the only ambiguity being a
componentwise positive scaling---a gauge symmetry (Theorem~\ref{ch:thm:uniqueness}).

\paragraph{3. A forced Hilbert-space geometry.}
The logarithmic factorisation $\log n_{mc}=v_c-u_m$ is not a rewriting:
it forces the definition of a canonical linear operator $\Phi$, and
from this operator every subsequent analytic construction follows
inexorably---the least-squares principle, the incidence Laplacian,
convexity modulo gauge, stability, the quotient geometry, and the UOD.
None of these is chosen; each is the unique consequence of the
logarithmic representation together with elementary Hilbert-space
theory.

\paragraph{4. The UOD as a canonical distance.}
The Universal Optimality Defect $\mathcal D(P,Q)$ is not another
divergence: it is the squared Euclidean distance in the quotient space
of logarithmic coordinates.  All smooth information functionals become
scalar fields on this geometry, and all inherit universal stability
bounds from two general theorems.

\paragraph{5. Universal Lipschitz and quadratic bounds.}
For any $C^1$ functional $S$ on a compact subset,
$|S(P)-S(Q)|\le L_S\sqrt{\mathcal D(P,Q)}$ (Theorem~\ref{ch:thm:14-1}).  For any
$C^2$ functional with a stationary point,
$|S(P)-S(P^\star)|\le C_S\,\mathcal D(P,P^\star)$ (Theorem~\ref{ch:thm:14-2}).
These bounds apply uniformly to Shannon entropy, R\'enyi entropy (all
$\alpha>0$), Tsallis entropy, min-entropy, KL divergence, $\chi^2$
divergence, Hellinger distance, $\alpha$-divergences, $\beta$-divergences,
Jensen--Shannon divergence, and R\'enyi divergence.

\paragraph{6. Lower bounds for convex and concave functionals.}
For geodesically convex functionals (KL, Jensen--Shannon),
$(\mu/2)\mathcal D(P,Q)\le S(P)-S(Q)\le (M/2)\mathcal D(P,Q)$
(Theorem~\ref{ch:thm:15-3}).  For geodesically concave functionals (R\'enyi
$\alpha<1$, min-entropy), the inequality is reversed.  This brackets
every major information functional between two multiples of the UOD.

\paragraph{7. A unified g-entropy framework.}
All classical entropies---Shannon, R\'enyi, Tsallis,
min-entropy---are special cases of a single functional
$H_g(P)=g(\sum_i g^{-1}(-\log P_i))$, whose Lipschitz constant is
given in closed form (Corollary~\ref{ch:cor:g-entropy-lipschitz}).  The bounds for each entropy are
consequences of one common structure.

\paragraph{8. Piecewise-smooth security measures.}
Rank-based measures (Bayes vulnerability, guessing entropy, max
leakage) are not globally smooth but become smooth on regularity
intervals where the posterior ranking is constant.  On each interval
the universal bounds apply, with explicit crossing times at interval
boundaries (Chapter~\ref{ch:sec:piecewise-smooth}).

\paragraph{9. Empirical validation of bound tightness.}
A computational discovery campaign on one-heavy distributions
$P_\varepsilon=(1-(n-1)\varepsilon,\varepsilon,\dots,\varepsilon)$
confirms that the Lipschitz bound is always satisfied but
$60$--$400\,000\times$ wider than the actual gap.  The quadratic bound
is $2$--$42\times$ wider than the exact R\'enyi identity.  The bounds
are structural---they show that all smooth functionals obey the same
geometric principle---not replacements for specialised inequalities.

\paragraph{10. The Three-Level Theorem.}
Every interior regular constrained critical point of the variational
problem $\max\mathcal D(P,U_n)$ subject to
$D_{\mathrm{KL}}(P\|U_n)=K$ collapses to at most three distinct
log-likelihood values, independently of $n$ (Theorem~\ref{ch:thm:29-3}).  The
critical points form a two-parameter family, and the transition from
one level (uniform) to two levels to three levels as $K$ increases is
a pitchfork bifurcation.

\paragraph{11. Randomised encoders and rational WIS.}
Every rational WIS is realisable by a randomised encoder with finite
seed set (Theorem~\ref{ch:thm:randomised-uod}).  Every real WIS weight is approximable with
error $O(1/|\mathcal S|)$.  The deterministic ACE are the integral
extremal points of the rational convex hull.

\paragraph{12. Eight concrete applications.}
The framework is demonstrated on AES S-box, browser fingerprinting,
geo-indistinguishability, differential privacy (both central and local
models), anonymous communication (mix networks), network routing, and
side-channel analysis.  Each follows the pattern
\textbf{problem} $\to$ \textbf{WIS} $\to$ \textbf{UOD}
$\to$ \textbf{interpretation}.

\section{Comparison with the State of the Art}\label{ch:sec:comparison-with-the-state-of-the-art}

\begin{center}
\begin{tabular}{lll}
\textbf{Aspect} & \textbf{Previous literature} & \textbf{This book} \\ \hline
Universally optimal encoders & ACE family only \cite{biondi2018apollonian} & Complete characterisation \\
WIS as invariant & Not identified & Intrinsic invariant \\
Geometry & Metric from divergence & Geometry from linearisation \\
Stability bounds & Separate per functional & Universal Lipschitz+quadratic \\
R\'enyi bounds & Exact via divergence \cite{vanerven2014renyi} & Independent UOD estimate \\
Security measures & Case-by-case & Unified piecewise-smooth theory \\
Lower bounds & Bregman only & All convex/concave functionals \\
Three-Level Structure & Does not exist & New variational result \\
Randomised encoders & Not studied & Rational approximation \\
\end{tabular}
\end{center}

The contribution is not a sharper inequality for a particular quantity,
but a structural unification: the same geometry---the WIS manifold with
its pullback metric and the UOD---simultaneously classifies ciphers,
bounds entropy, estimates security, and reveals the variational
collapse of the uniform observer divergence.

\section{Organisation of the Book}\label{ch:sec:organisation-of-the-book}

The book is organised in five parts, preceded by a chapter of
motivating examples (Chapter~\ref{ch:sec:motivating-examples}) and the present introduction
(Chapter~\ref{ch:sec:introduction}).

\paragraph{Part~\ref{ch:part:wis}---Weighted Incidence Structures (Chapters~\ref{ch:sec:background}--\ref{ch:sec:representation}).}
Develops the probabilistic setting, multiplicity matrices, the
cycle-factorisation criterion, weighted incidence structures, the
classification theorem, and gauge uniqueness.

\paragraph{Part~\ref{ch:part:hilbert}---Hilbert Geometry and the Universal Optimality
Defect (Chapters~\ref{ch:sec:graphs}--\ref{ch:sec:hilbert-geometry}).}
Constructs the logarithmic linearisation, the mismatch operator, the
least-squares variational principle, convexity, stability, and the
Universal Optimality Defect, with worked examples.

\paragraph{Part~\ref{ch:part:posterior}---Geometric Foundations of the Posterior Manifold
(Chapters~\ref{ch:sec:the-posterior-manifold}--\ref{ch:sec:stability-of-the-universal-optimality-defect}).}
Constructs a Riemannian manifold on the interior of the probability
simplex via the pullback of the Euclidean metric.  The UOD becomes the
squared distance; its stability, local quadratic approximation, and
relation with Fisher geometry are developed.

\paragraph{Part~\ref{ch:part:functionals}---Information Functionals and Geometric Bounds
(Chapters~\ref{ch:sec:entropy-bounds}--\ref{ch:sec:universal-security-bounds}).}
Applies the manifold geometry to information-theoretic and security
functionals.  Universal Lipschitz and quadratic bounds are established
for all classical entropies, divergences, and security measures.
Bregman divergences are analysed through the Average Hessian Operator.
Lower bounds for convex and concave functionals complete the bracketing.

\paragraph{Part~\ref{ch:part:variational}---Variational Structure and Comparative Geometry
(Chapters~\ref{ch:sec:analytic-uod}--\ref{ch:sec:open}).}
Develops the algebraic and differential theory of the UOD, establishes
local and global comparisons with classical statistical divergences,
and proves the Three-Level Theorem.  The final four chapters of this part
present conclusions, a discussion of related work, eight extended
applications (AES S-box, browser fingerprinting, geo-indistinguishability,
differential privacy, mix networks, routing, side channels, local DP),
and open problems including randomised encoders, quantum WIS,
computational complexity, and gradient flows.

\paragraph{Relation to the ACE literature.}
The present work is not an extension of the Apollonian Cell Encoder
construction \cite{biondi2018apollonian}.  ACE provides one explicit
family of universally optimal encoders.  This book asks the converse
structural question: which deterministic encoders are universally
optimal, and what mathematical structure do they all share?  The answer
is the weighted incidence structure and its forced Hilbert-space
geometry, which subsume the ACE construction as a special case.
\endinput